% Core chapter-5 teaching sequence, included by lecture_notes.tex.
\section{Pretraining: loss and evaluation, chapter 5.1}
\label{sec:pretraining-loss}

\begin{studybox}{What pretraining changes}
Start with randomly initialized model weights and a fixed tokenizer. The model
learns to predict the next token from observed prefixes of ordinary text. It
does not receive instruction/answer pairs, preference labels, or a reward model.
The targets come from shifting the text itself. A trained tokenizer and a
trained language model are separate artifacts.
\end{studybox}

% Original course diagram. All text, mathematics, and connectors remain vector.
\begin{figure}[htbp]
\centering
\begin{tikzpicture}[
  font=\small,
  flow/.style={-{Latex[length=2mm]},thick,draw=navy},
  target/.style={-{Latex[length=2mm]},draw=navy,dashed},
  update/.style={-{Latex[length=2mm]},thick,draw=coral!85!black},
  box/.style={draw=navy,fill=navy!5,rounded corners=2pt,
    align=center,minimum height=0.88cm,inner sep=4pt},
  train/.style={box,draw=teal!80!black,fill=teal!8},
  check/.style={box,draw=navy,fill=codebg}
]
\node[box,text width=1.5cm] (corpus) at (1,0) {Text\\corpus};
\node[box,text width=1.8cm] (split) at (3.25,0) {Split text\\before windows};
\node[train,text width=2.1cm] (trainwin) at (6,1.3)
  {Training\\windows $(x,y)$};
\node[check,text width=2.1cm] (valwin) at (6,-1.3)
  {Validation\\windows $(x,y)$};
\node[train,text width=1.9cm] (trainmodel) at (9,1.3)
  {GPT $\theta$\\logits $z$};
\node[check,text width=1.9cm,minimum height=1.08cm] (valmodel) at (9,-1.3)
  {GPT $\theta$\\no gradients\\logits $z$};
\node[train,text width=1.4cm] (trainloss) at (11.7,1.3)
  {$\mathrm{CE}(z,y)$\\training loss};
\node[check,text width=1.4cm] (valloss) at (11.7,-1.3)
  {$\mathrm{CE}(z,y)$\\held-out loss};
\node[train,text width=1.55cm,draw=coral!85!black,fill=coral!8]
  (optim) at (14.05,1.3) {Backward\\AdamW step};
\node[align=center,text width=1.65cm,text=navy] (report) at (14.05,-1.3)
  {Report and\\compare};

\draw[flow] (corpus) -- (split);
\draw[flow] (split.east) -- ++(0.35,0) |- (trainwin.west);
\draw[flow] (split.east) -- ++(0.35,0) |- (valwin.west);
\draw[flow] (trainwin) -- node[above] {$x$} (trainmodel);
\draw[flow] (valwin) -- node[above] {$x$} (valmodel);
\draw[flow] (trainmodel) -- (trainloss);
\draw[flow] (valmodel) -- (valloss);
\draw[flow] (trainloss) -- (optim);
\draw[flow] (valloss) -- (report);

% Shifted target IDs go directly to the loss, not into the decoder input.
\draw[target] (trainwin.south) -- (6,0.2)
  -- node[below] {targets $y$} (11.7,0.2) -- (trainloss.south);
\draw[target] (valwin.south) -- (6,-2.4)
  -- node[below] {targets $y$} (11.7,-2.4) -- (valloss.south);
\draw[update] (optim.north) -- (14.05,2.6)
  -- node[above] {update $\theta$; repeat with the next batch} (9,2.6)
  -- (trainmodel.north);
\end{tikzpicture}
\caption{The pretraining data and optimization paths. Split the corpus before
forming overlapping windows, tokenize each split with the same fixed tokenizer,
and shift each window by one position to obtain input IDs $x$ and target IDs $y$.
Cross-entropy (CE) compares the logits with the targets. Only the training branch
computes gradients and updates parameters. Validation uses the same current
parameters in evaluation mode, with dropout disabled and no gradient recording.
If the tokenizer is learned from the course corpus, fit it only on the training
split.}
\label{fig:pretraining-flow}
\end{figure}


\subsection{From a sequence probability to a trainable loss}
The probability chain rule is exact:
\[
p(x_1,\ldots,x_T)=\prod_{t=1}^{T}p(x_t\mid x_{<t}).
\]
The model approximates these conditional distributions using a finite context
and shared parameters $\theta$. A training input window $(x_1,\ldots,x_T)$ has
labels $(x_2,\ldots,x_{T+1})$. For its observed next tokens, maximizing the
product of conditional probabilities is equivalent to maximizing their summed
logarithms, hence to minimizing negative log-likelihood:
\begin{align*}
\arg\max_\theta\prod_{b,t}p_\theta(y_{bt}\mid x_{b,\le t})
&=\arg\max_\theta\sum_{b,t}\log p_\theta(y_{bt}\mid x_{b,\le t}),\\
\mathcal L_{\rm LM}(\theta)
&=-\frac{1}{BT}\sum_{b=1}^{B}\sum_{t=1}^{T}
\log p_\theta(y_{bt}\mid x_{b,\le t}).
\end{align*}
In these full-length pretraining windows every target contributes. Later, padding
or assistant-only training replaces $BT$ by the number of included targets and
omits masked terms. Such loss masks are different from the attention mask.

Training uses \emph{teacher forcing}: each position sees the observed earlier
tokens even if the model would have generated something else. All positions
can be evaluated in parallel because causal attention blocks future inputs.
Generation instead appends the model's own selected tokens. Errors may then
change the prefixes encountered later in the rollout.

\subsection{Logits, target probability, and cross-entropy}
Let $z_{btj}$ be the raw logit for vocabulary item $j$. For one target ID $y$,
\[
p_j=\frac{e^{z_j}}{\sum_{k=0}^{V-1}e^{z_k}},\qquad
\ell(z,y)=-\log p_y
=\log\left(\sum_{k=0}^{V-1}e^{z_k}\right)-z_y.
\]
The last expression connects the next-token objective to cross-entropy against
the one-hot target distribution. Implement the combined operation with
\texttt{cross\_entropy}; do not explicitly exponentiate large logits, take a
logarithm of a rounded probability, or pass softmax probabilities to that API.

For example, if the model assigns its three correct next tokens probabilities
$1/2$, $1/4$, and $1/8$, their mean loss is
\[
-\tfrac13\left(\log\tfrac12+\log\tfrac14+\log\tfrac18\right)
=\log4\approx1.3863.
\]
This is a constructed numerical example, not a measured training result.
For a single logit vector the derivative is
\[
\frac{\partial\ell}{\partial z_j}=p_j-\mathbf1\{j=y\}.
\]
Thus gradient descent increases the correct token's score relative to the
others, while backpropagation adjusts every upstream component that contributed
to these logits. Autograd performs this chain rule; it does not choose targets,
the loss, or the optimizer for us.

\begin{lstlisting}[language=Python]
import torch
from torch.nn import functional as F

def loss_batch(model, x, y):
    logits = model(x).logits  # book: logits = model(x)
    # logits: (B, T, V), y: (B, T), integer next-token IDs
    return F.cross_entropy(logits.flatten(0, 1), y.flatten())
\end{lstlisting}
The flattening creates $BT$ classification events over $V$ classes. It does
not mix time steps within a target and does not shift labels. A second shift
here would train the wrong prediction task.

\subsection{What perplexity measures}
For natural-log loss, $\operatorname{PPL}=\exp(\mathcal L)$. The example above
has perplexity $4$, the inverse geometric mean of the three correct-token
probabilities. A uniform predictor has loss $\log V$ and perplexity $V$.
A randomly initialized network need not be exactly uniform.

Perplexity is not accuracy and does not count the literal number of plausible
words. It depends on token boundaries, the held-out text, the context policy,
and which targets are scored. Compare it only under the same evaluation
protocol. Mean token loss should be exponentiated once; averaging per-batch
perplexities is a different quantity.

\subsection{The book's small corpus and split}
Raschka's chapter-5 example uses Edith Wharton's short story \emph{The Verdict}.
The exact course-reading download source is the pinned
\booklink{ch02/01_main-chapter-code/the-verdict.txt}{repository text file};
the raw endpoint is linked in the source records at the end of these notes.
The demonstration divides the raw text at 90\% of its character length,
tokenizes the two pieces separately with the fixed GPT-2 tokenizer, and then
constructs windows. This prevents overlapping windows from crossing the split,
but the two sets still come from the same story. It is not a held-out-author
or held-out-book generalization experiment.

The chapter-5 notebook uses context length and stride $256$, batch size $2$,
training shuffle enabled, and incomplete training batches dropped. Validation
does not shuffle and does not drop the last batch. An epoch is one pass through
the resulting training loader, not one pass through every possible overlapping
window or every possible prefix length.

\subsection{Evaluation must change both mode and gradient tracking}
\texttt{model.eval()} disables dropout; it does not disable autograd.
\texttt{torch.no\_grad()} disables graph recording; it does not disable dropout.
Use both to measure loss with fixed model parameters and deterministic network
layers. Restore the previous mode afterward.

The book's \texttt{calc\_loss\_loader} averages batch means over at most
\texttt{eval\_iter} batches. That is an exact token mean when batches contain
the same number of scored targets. The following original course adaptation
weights by token count so an incomplete final batch is handled correctly.
It assumes dense pretraining labels, with no padding or ignored targets.

\begin{lstlisting}[language=Python]
@torch.no_grad()
def evaluate(model, loader, device, max_batches=None):
    was_training = model.training
    model.eval()
    total_nll, total_tokens = 0.0, 0
    try:
        for i, (x, y) in enumerate(loader):
            if max_batches is not None and i >= max_batches:
                break
            x, y = x.to(device), y.to(device)
            logits = model(x).logits  # book: model(x)
            total_nll += F.cross_entropy(
                logits.flatten(0, 1), y.flatten(),
                reduction="sum").item()
            total_tokens += y.numel()
    finally:
        model.train(was_training)
    if total_tokens == 0:
        raise ValueError("No validation targets were scored")
    return total_nll / total_tokens
\end{lstlisting}

When \texttt{max\_batches=5}, five is a maximum, not a guarantee. In the book's
short-story example the validation loader is smaller than that. Partial-loader
or sampled-window evaluation is a useful progress estimate, but must not be
reported as an exhaustive corpus evaluation. Reusing validation to choose
hyperparameters makes it a model-selection set; reserve a separate test set for
a final claim.

\section{Pretraining: the first training loop, chapter 5.2}
\label{sec:basic-loop}

\subsection{Start with one device and one optimizer step per batch}
The first loop should expose the essential computation before adding a schedule,
mixed precision, gradient accumulation, or multiple GPUs:
\begin{enumerate}
\item Obtain integer inputs $X$ and already-shifted targets $Y$.
\item Compute raw logits and the mean cross-entropy.
\item Clear old gradients, backpropagate this batch's loss, and update parameters.
\item Periodically evaluate and record the number of processed token positions.
\item Inspect continuations from a fixed prompt after each epoch.
\end{enumerate}
Clearing gradients is necessary because PyTorch accumulates them by default.
\texttt{backward()} computes gradients but does not change parameter values;
\texttt{optimizer.step()} makes the update. For plain gradient descent this is
$\theta\leftarrow\theta-\eta\nabla\mathcal L$. AdamW additionally tracks first
and second gradient moments and applies decoupled weight decay.

The \emph{chapter notebook} uses AdamW with learning rate $4\times10^{-4}$,
weight decay $0.1$, ten epochs, and evaluation every five updates using at most
five batches per loader. It seeds PyTorch before constructing the model.
The compact \texttt{gpt\_train.py} summary uses different defaults
($5\times10^{-4}$ and one evaluation batch); do not mix those settings and call
the result the notebook experiment. The original short-story setup has a large
model relative to its data and intentionally exposes overfitting.

\begin{lstlisting}[language=Python]
def train_basic(model, train_loader, val_loader, device,
                epochs=10, eval_every=5):
    model.to(device)
    optimizer = torch.optim.AdamW(
        model.parameters(), lr=4e-4, weight_decay=0.1)
    history, tokens_seen, updates = [], 0, 0
    for epoch in range(epochs):
        model.train()
        for x, y in train_loader:
            x, y = x.to(device), y.to(device)
            optimizer.zero_grad(set_to_none=True)
            loss = loss_batch(model, x, y)
            if not torch.isfinite(loss):
                raise RuntimeError("Non-finite training loss")
            loss.backward()
            optimizer.step()
            tokens_seen += y.numel()
            updates += 1
            if updates == 1 or updates % eval_every == 0:
                train_nll = evaluate(model, train_loader,
                                     device, max_batches=5)
                val_nll = evaluate(model, val_loader,
                                   device, max_batches=5)
                history.append((updates, tokens_seen,
                                train_nll, val_nll))
        # Inspect a fixed-prompt continuation here; see below.
    return history, optimizer
\end{lstlisting}
This is an original condensed loop with the course \texttt{GPTOutput} interface,
not a verbatim copy of the book. It uses one-based update numbers; the book's
zero-based \texttt{global\_step=0} already follows its first optimizer update.
For a fresh run, set the seed \emph{before} creating the model, loaders, and
optimizer. An evaluation before the loop is useful as a genuine untrained
baseline. Calling the first post-update measurement ``initial loss'' would
confuse these two points in time.

\subsection{Reading training curves without overclaiming}
Compare train and validation loss with dropout disabled and the same scoring
convention. Falling train loss with rising held-out loss is evidence of
overfitting under that protocol. Similar losses can be entirely healthy; they
do not alone prove leakage. A nearly constant loss may reflect an incorrect
target shift, missing updates, unsuitable learning rate, or too little training.
Inspect these mechanisms before scaling the model.

The course's \emph{production-style driver} logs the just-used training batches
with dropout before an update, and evaluates sampled validation batches without
dropout after it. Those two logged values are not a clean matched estimate of
the generalization gap. Its validation batches also change over time. Use a
fixed held-out scoring protocol for comparisons that need finer precision.

\begin{studybox}{A useful small experiment}
Overfit one small batch to test whether the network and optimizer can reduce
its loss. Then evaluate on text never used for updates. Success on the repeated
batch is a correctness diagnostic, not generalization evidence. Compare several
fixed-prompt continuations, not only the most fluent one. Keep the tokenizer,
prompt, context length, and decoding settings fixed across checkpoints.
\end{studybox}

\section{From a checkpoint to text: chapters 5.3--5.4}
\label{sec:book-generation}

\subsection{Greedy decoding, temperature, and top-k}
For a prompt of length $t$, run the last $\min(t,C)$ IDs through the decoder.
The vector \texttt{logits[:, -1, :]} predicts the next ID. Select that ID, append
it, and repeat. Earlier output positions are useful for parallel training but
do not predict the next token after the entire current prompt.

Greedy decoding chooses $\arg\max_j z_j$. For $\tau>0$, temperature sampling uses
\[
p_j=[\softmax(z/\tau)]_j,\qquad
x_{t+1}\sim\operatorname{Categorical}(p).
\]
Lower temperature concentrates probability and higher temperature flattens it.
Changing a positive temperature cannot change the greedy argmax ordering.
It can change a sampled sequence, but does not guarantee a different realized
draw. Top-$k$ keeps the largest $k$ scores and masks the rest before normalizing.
At $k=1$ this is effectively greedy when the maximum is unique.

\begin{lstlisting}[language=Python]
@torch.no_grad()
def sample_next(model, ids, context_size, vocab_size,
                temperature=1.0, top_k=40):
    if temperature <= 0:
        raise ValueError("Use argmax for greedy decoding")
    if not 1 <= top_k <= vocab_size:
        raise ValueError("top_k must be within the vocabulary")
    if vocab_size > model.config.vocab_size:
        raise ValueError("Tokenizer is larger than the model")
    was_training = model.training
    model.eval()
    try:
        z = model(ids[:, -context_size:]).logits[:, -1, :]
        z = z[:, :vocab_size] / temperature
        values, indices = torch.topk(z, top_k, dim=-1)
        filtered = torch.full_like(z, -torch.inf)
        filtered.scatter_(1, indices, values)
        return torch.multinomial(filtered.softmax(-1), 1)
    finally:
        model.train(was_training)
\end{lstlisting}
This course-interface snippet selects one token; call it repeatedly and
concatenate its return value along dimension 1. It intentionally recomputes the
cropped context and has no key/value cache. It assumes a nonempty prompt,
$1\le\texttt{context\_size}\le C$, valid input IDs, and finite logits. The
full CLI additionally handles checkpoint/tokenizer validation and stopping.
Released GPT-2 token IDs cannot be substituted for the course tokenizer's IDs.

\subsection{Saving weights versus resuming training}
Chapter 5.4 first saves and reloads \texttt{model.state\_dict()}, then includes
optimizer state for continued training. Reconstruct the same model configuration
before loading its weights. For inference, select evaluation mode and pair it
with the exact tokenizer. Loading only weights and creating a new optimizer is
a fresh optimizer trajectory, even if the parameter values are unchanged.

A course checkpoint records configuration, model state, optimizer state, update
number, and metadata including the tokenizer identity. It does \emph{not} yet
provide exact restart: the current pretraining CLI has no resume option and its
payload omits RNG, batch-generator, and GradScaler states. An exact-resume
implementation would also restore the schedule position and data-order state.
Only load checkpoint files from a trusted source because general serialized
PyTorch payloads can execute code when deserialized.

\subsection{The boundary of ``from scratch''}
Chapter 5.5 separately demonstrates loading released OpenAI GPT-2 weights.
That is a useful architecture-compatibility exercise, not a result obtained by
training on \emph{The Verdict}. The course's pretraining exercise starts from
random weights and does not require that step. Similarly, useful instruction
following is a later objective rather than an automatic consequence of a
working next-token training loop.

\begin{studybox}{Pretraining checkpoint questions}
Can you explain every dimension from input IDs to loss? Why must masking occur
before softmax? Why does \texttt{eval()} not replace \texttt{no\_grad()}?
What is counted by ``tokens seen''? What would make two perplexities comparable?
Which states would you restore to resume training exactly? Answer these before
adding the larger-corpus and hardware extensions in the next sections.
\end{studybox}
